\documentclass[aspectratio=169,11pt]{beamer}

% Compile with XeLaTeX
\usepackage{fontspec}
\setsansfont{DejaVu Sans}
\setmonofont{DejaVu Sans Mono}
\usepackage{polyglossia}
\setdefaultlanguage{english}
\usepackage{booktabs}
\usepackage{array}
\usepackage{tabularx}
\usepackage{adjustbox}
\usepackage{listings}
\usepackage{xcolor}
\usepackage{ragged2e}
\usepackage{graphicx}
\usepackage{hyperref}

% Graphics
\graphicspath{{images/}{./}{/mnt/data/images/}}
\hypersetup{unicode=true}

\usetheme{Madrid}
\usecolortheme{seahorse}
\setbeamertemplate{navigation symbols}{}
\setbeamertemplate{footline}[frame number]
\setbeamertemplate{itemize items}[circle]
\setbeamertemplate{enumerate items}[default]

\definecolor{codebg}{RGB}{248,248,248}
\definecolor{codekw}{RGB}{0,72,160}
\definecolor{codestr}{RGB}{150,60,30}
\definecolor{accent}{RGB}{0,90,150}

\lstset{
  basicstyle=\ttfamily\scriptsize,
  keywordstyle=\color{codekw}\bfseries,
  stringstyle=\color{codestr},
  backgroundcolor=\color{codebg},
  frame=single,
  breaklines=true,
  columns=fullflexible,
  showstringspaces=false,
  keepspaces=true,
  language=SQL
}

\newcommand{\keyword}[1]{\textcolor{accent}{\textbf{#1}}}
\newcommand{\kw}[1]{\keyword{#1}}
\newcommand{\tightlist}{\setlength{\itemsep}{2pt}\setlength{\parskip}{0pt}}
\newcommand{\smallitem}{\setlength{\itemsep}{2pt}\setlength{\parskip}{0pt}}
\newcommand{\qoption}[2]{\item[#1.] #2}
\newcommand{\safeimage}[3]{%
\IfFileExists{#1}{\includegraphics[width=#2,height=0.62\textheight,keepaspectratio]{#1}}{%
\fbox{\parbox[c][0.40\textheight][c]{0.78\textwidth}{\centering #3\\[2mm]{\scriptsize #1}}}}}

\title[Choosing a DBMS]{Choosing the Right Database Management System}
\subtitle{Selection Criteria, Database Types, and Operational Models}
\author{Lecture Notes}
\institute{Course: Databases}
\date{Updated: 31/05/2026}

\AtBeginSection[]{
  \begin{frame}{Section Outline}
    \tableofcontents[currentsection,subsectionstyle=show/show/hide]
  \end{frame}
}

\begin{document}

\begin{frame}
  \titlepage
\end{frame}

% \begin{frame}{References}
% \begin{itemize}
%   \item GeeksforGeeks: \emph{How to Choose the Database}.
%   \item The lecture content has been reorganized into an instructional format.
% \end{itemize}
% \end{frame}

\begin{frame}{Learning Objectives}
After completing this lesson, learners will be able to:
\begin{enumerate}
  \smallitem
  \item Understand why choosing a database management system is an important architectural decision.
  \item Distinguish common database types: relational, document, key-value, time-series, columnar, and graph databases.
  \item Identify the main criteria for selecting a DBMS for a software system.
  \item Understand when to use a self-managed database and when to use a managed database service.
  \item Apply the knowledge to choose a suitable database type for a practical problem.
  \item Complete review questions and application exercises on DBMS selection.
\end{enumerate}
\end{frame}

\begin{frame}[allowframebreaks]{Lecture Structure}
\small
\tableofcontents
\end{frame}

\section{Overview}

\begin{frame}{Why Does Choosing the Right DBMS Matter?}
In any software system that uses data, selecting the right \kw{database management system} is a very important decision.

\vspace{2mm}
A suitable choice can help the system:
\begin{itemize}
  \smallitem
  \item Store data safely and consistently.
  \item Query data faster.
  \item Scale more easily as the number of users grows.
  \item Reduce operational costs.
  \item Become easier to maintain and evolve in the future.
\end{itemize}
\end{frame}

\begin{frame}{Consequences of Choosing the Wrong DBMS}
If the wrong database type is selected, the system may face many problems:
\begin{itemize}
  \smallitem
  \item Slow queries.
  \item Difficulty scaling.
  \item Inconsistent data.
  \item High operational costs.
  \item Difficulty changing the data structure as the system evolves.
\end{itemize}

\begin{block}{Key Message}
Do not choose a database simply because it is popular. The choice should depend on the nature of the data, query patterns, performance, consistency requirements, scalability, and operational capability.
\end{block}
\end{frame}

\begin{frame}{Factors to Consider When Choosing a Database}
The decision should be based on:
\begin{columns}[T]
\begin{column}{0.48\textwidth}
\begin{itemize}
  \smallitem
  \item The data type used by the system.
  \item How the data is queried.
  \item Performance requirements.
  \item Consistency requirements.
\end{itemize}
\end{column}
\begin{column}{0.48\textwidth}
\begin{itemize}
  \smallitem
  \item Scalability requirements.
  \item Operational capability of the technical team.
  \item Cost.
  \item Long-term system development direction.
\end{itemize}
\end{column}
\end{columns}
\end{frame}

\begin{frame}{Figure 1. Choosing the Right Database}
\centering
\safeimage{images/choose-your-right-database.png}{0.82\textwidth}{Choose your right database}

\vspace{2mm}
{\small Overview of choosing an appropriate database.}
\end{frame}

\begin{frame}{Quick Quiz: Overview}
\small
\textbf{Question 1.} Why should we not choose a DBMS only because it is popular?
\begin{enumerate}
  \qoption{A}{Because popular DBMSs are always free}
  \qoption{B}{Because DBMS selection should depend on data, queries, performance, and operations}
  \qoption{C}{Because popular DBMSs do not support security}
  \qoption{D}{Because every system should use only one database type}
\end{enumerate}

\textbf{Question 2.} What problem may occur when the DBMS choice is unsuitable?
\begin{enumerate}
  \qoption{A}{Slow queries and difficulty scaling}
  \qoption{B}{Higher monitor resolution}
  \qoption{C}{Fewer tables required in every system}
  \qoption{D}{Automatic removal of the need for backups}
\end{enumerate}
\end{frame}

\section{Basic Concepts}

\begin{frame}{What Is a DBMS?}
\kw{DBMS} stands for \kw{Database Management System}.

\vspace{2mm}
A DBMS is software used to:
\begin{columns}[T]
\begin{column}{0.48\textwidth}
\begin{itemize}
  \smallitem
  \item Create databases.
  \item Store data.
  \item Query data.
  \item Update data.
\end{itemize}
\end{column}
\begin{column}{0.48\textwidth}
\begin{itemize}
  \smallitem
  \item Delete data.
  \item Ensure data security.
  \item Ensure data integrity.
  \item Manage access and transactions.
\end{itemize}
\end{column}
\end{columns}
\end{frame}

\begin{frame}{Examples of DBMSs}
\begin{columns}[T]
\begin{column}{0.48\textwidth}
\begin{itemize}
  \smallitem
  \item MySQL
  \item PostgreSQL
  \item SQL Server
  \item Oracle Database
\end{itemize}
\end{column}
\begin{column}{0.48\textwidth}
\begin{itemize}
  \smallitem
  \item MongoDB
  \item Redis
  \item Neo4j
  \item Snowflake
\end{itemize}
\end{column}
\end{columns}
\end{frame}

\begin{frame}{Quick Quiz: Basic Concepts}
\small
\textbf{Question 1.} What does DBMS stand for?
\begin{enumerate}
  \qoption{A}{Database Management System}
  \qoption{B}{Data Backup Monitoring Service}
  \qoption{C}{Distributed Binary Model Server}
  \qoption{D}{Database Mapping Syntax}
\end{enumerate}

\textbf{Question 2.} Which of the following is a function of a DBMS?
\begin{enumerate}
  \qoption{A}{Creating, storing, querying, updating, and protecting data}
  \qoption{B}{Designing graphical user interfaces}
  \qoption{C}{Compiling application source code}
  \qoption{D}{Completely replacing the operating system}
\end{enumerate}
\end{frame}

\section{Criteria for Selecting a DBMS}

\begin{frame}{Overview of the Main Criteria}
When choosing a database for a system, consider:
\begin{enumerate}
  \smallitem
  \item Scalability.
  \item Data consistency.
  \item Query patterns and workload.
  \item Operational complexity.
  \item Recoverability and availability.
  \item Performance requirements.
\end{enumerate}
\end{frame}

\begin{frame}{Criterion 1: Scalability}
Scalability indicates whether the system can handle increasing data volume or a growing number of users.

\vspace{2mm}
There are two main directions for scaling:
\begin{block}{Vertical Scaling}
Increasing the resources of an existing server.
\end{block}
\begin{block}{Horizontal Scaling}
Adding more servers to the system to distribute data or load.
\end{block}
\end{frame}

\begin{frame}{Vertical Scaling}
Vertical scaling means increasing the resources of an existing server.

\vspace{2mm}
Examples:
\begin{itemize}
  \smallitem
  \item Adding more CPU.
  \item Adding more RAM.
  \item Increasing disk capacity.
  \item Using a more powerful server.
\end{itemize}

\begin{block}{Remember}
Relational databases such as MySQL, PostgreSQL, and SQL Server are often suitable for vertical scaling.
\end{block}
\end{frame}

\begin{frame}{Horizontal Scaling}
Horizontal scaling means adding more servers to the system.

\vspace{2mm}
Examples:
\begin{itemize}
  \smallitem
  \item Splitting data across multiple machines.
  \item Distributing query load.
  \item Using a server cluster.
\end{itemize}

\begin{block}{Remember}
NoSQL databases are often designed to support horizontal scaling well.
\end{block}
\end{frame}

\begin{frame}{Criterion 2: Data Consistency}
Data consistency indicates whether data is accurate and synchronized at all times.

\begin{exampleblock}{Banking Example}
When account A transfers money to account B, the amount must be deducted from A and added to B. These two operations must occur together.
\end{exampleblock}

If money is deducted from A but not added to B, the data becomes incorrect.
\end{frame}

\begin{frame}{Strong Consistency and Eventual Consistency}
\begin{block}{Strong Consistency}
Financial, accounting, and banking systems usually require strong consistency and ACID-compliant transactions.
\end{block}

\begin{block}{Eventual Consistency}
Some systems can tolerate data becoming consistent after a short delay, such as social networks, recommendation systems, logging systems, or view-count analytics.
\end{block}
\end{frame}

\begin{frame}{Criterion 3: Query Patterns and Workload}
Before choosing a DBMS, we need to understand what the system mainly does with its data.

\vspace{2mm}
Common workload patterns include:
\begin{itemize}
  \smallitem
  \item Read-heavy workloads.
  \item Write-heavy workloads.
  \item Balanced read and write workloads.
  \item Complex queries.
  \item Simple queries requiring very high speed.
\end{itemize}
\end{frame}

\begin{frame}{Query Patterns and Suitable Database Types}
\scriptsize
\begin{adjustbox}{max width=\textwidth}
\begin{tabularx}{\textwidth}{>{\raggedright\arraybackslash}p{0.44\textwidth}>{\raggedright\arraybackslash}X}
\toprule
\textbf{System Characteristic} & \textbf{Suitable Database Type} \\
\midrule
Requires many complex JOIN operations & Relational database \\
Requires reporting, statistics, and analytics & Relational database or columnar database \\
Requires very fast key-based access & Key-value store \\
Uses flexible JSON data & Document database \\
Uses sensor data over time & Time-series database \\
Contains many complex relationships & Graph database \\
\bottomrule
\end{tabularx}
\end{adjustbox}
\end{frame}

\begin{frame}{Criterion 4: Operational Complexity}
A database system not only needs to be well designed; it also needs to be operated well.

\vspace{2mm}
Operational tasks include:
\begin{columns}[T]
\begin{column}{0.48\textwidth}
\begin{itemize}
  \smallitem
  \item Backing up data.
  \item Recovering data.
  \item Upgrading versions.
  \item Monitoring performance.
\end{itemize}
\end{column}
\begin{column}{0.48\textwidth}
\begin{itemize}
  \smallitem
  \item Managing user permissions.
  \item Encrypting data.
  \item Setting up replication.
  \item Handling incidents.
\end{itemize}
\end{column}
\end{columns}

\vspace{1mm}
If the technical team has limited database administration experience, a \kw{managed database} should be considered.
\end{frame}

\begin{frame}{Criterion 5: Recoverability and Availability}
Some systems must run continuously:
\begin{itemize}
  \smallitem
  \item Digital banking.
  \item E-commerce.
  \item Ticket booking systems.
  \item Hospital management systems.
  \item Real-time logistics systems.
\end{itemize}

Such systems need automatic backups, automatic failover, data replication, multi-region support, and disaster recovery capability.
\end{frame}

\begin{frame}{Criterion 6: Performance Requirements}
Performance can be evaluated through:
\begin{itemize}
  \smallitem
  \item Low latency.
  \item High throughput.
  \item Data write speed.
  \item Data read speed.
  \item Ability to process concurrent queries.
\end{itemize}

\begin{exampleblock}{Examples}
Cache systems require very low latency; IoT systems need very high write speed; BI systems require large analytical queries; transaction systems require consistency and reliability.
\end{exampleblock}
\end{frame}

\begin{frame}{Quick Quiz: DBMS Selection Criteria}
\small
\textbf{Question 1.} What does horizontal scaling mean?
\begin{enumerate}
  \qoption{A}{Adding more CPU and RAM to an existing server}
  \qoption{B}{Adding more servers to distribute load or data}
  \qoption{C}{Reducing the number of system users}
  \qoption{D}{Moving all data into text files}
\end{enumerate}

\textbf{Question 2.} What property is usually required in a banking system?
\begin{enumerate}
  \qoption{A}{Strong consistency and ACID transactions}
  \qoption{B}{Eventual consistency for all transactions only}
  \qoption{C}{No need for data integrity control}
  \qoption{D}{No need for backups}
\end{enumerate}
\end{frame}

\section{Common Database Types}

\begin{frame}{Figure 2. Common Database Types}
\centering
\safeimage{images/types-of-databases.png}{0.82\textwidth}{Types of databases}

\vspace{2mm}
{\small Common database types.}
\end{frame}

\begin{frame}{Database Types Covered in This Lesson}
\begin{columns}[T]
\begin{column}{0.48\textwidth}
\begin{itemize}
  \smallitem
  \item Relational database.
  \item Document database.
  \item Key-value store.
\end{itemize}
\end{column}
\begin{column}{0.48\textwidth}
\begin{itemize}
  \smallitem
  \item Time-series database.
  \item Columnar database.
  \item Graph database.
\end{itemize}
\end{column}
\end{columns}
\end{frame}

\begin{frame}{Quick Quiz: Common Database Types}
\small
\textbf{Question 1.} Tabular data with primary keys and foreign keys is usually suitable for which database type?
\begin{enumerate}
  \qoption{A}{Relational database}
  \qoption{B}{Time-series database}
  \qoption{C}{Graph database}
  \qoption{D}{Key-value store}
\end{enumerate}

\textbf{Question 2.} Flexible JSON data is usually suitable for which database type?
\begin{enumerate}
  \qoption{A}{Document database}
  \qoption{B}{Columnar database}
  \qoption{C}{Relational database in every case}
  \qoption{D}{Time-series database}
\end{enumerate}
\end{frame}

\subsection{Relational Database}

\begin{frame}{Relational Database: Concept}
A \kw{Relational Database Management System}, or \kw{RDBMS}, stores data in tables.

\vspace{2mm}
Each table includes:
\begin{itemize}
  \smallitem
  \item Rows.
  \item Columns.
  \item Primary keys.
  \item Foreign keys.
  \item Data constraints.
\end{itemize}

RDBMSs commonly use SQL to query data.
\end{frame}

\begin{frame}{Example: Students Table}
\scriptsize
\begin{center}
\begin{tabular}{lll}
\toprule
\textbf{student\_id} & \textbf{full\_name} & \textbf{major} \\
\midrule
1 & Nguyen An & Data Science \\
2 & Tran Binh & Logistics \\
3 & Le Chi & Information Systems \\
\bottomrule
\end{tabular}
\end{center}

\vspace{2mm}
Example SQL query:
\begin{block}{SQL}
\ttfamily SELECT full\_name, major FROM Students WHERE major = 'Data Science';
\end{block}
\end{frame}

\begin{frame}{When Should We Use an RDBMS?}
Use an RDBMS when:
\begin{itemize}
  \smallitem
  \item The data has a clear structure.
  \item Data integrity must be enforced.
  \item ACID transactions are required.
  \item JOIN operations across multiple tables are needed.
  \item Reporting and complex queries are required.
  \item The system involves finance, accounting, order management, or student management.
\end{itemize}
\end{frame}

\begin{frame}{Examples Suitable for RDBMSs}
\begin{itemize}
  \smallitem
  \item Banking systems.
  \item Student management systems.
  \item Sales systems.
  \item Inventory management systems.
  \item Accounting systems.
\end{itemize}
\end{frame}

\begin{frame}{When Should We Use a Managed RDBMS?}
Use a managed RDBMS service when:
\begin{itemize}
  \smallitem
  \item High availability is required.
  \item Automatic backups are desired.
  \item Point-in-time recovery is needed.
  \item The team does not want to manage version upgrades manually.
  \item Built-in security, encryption, and monitoring are required.
  \item The team wants to reduce operational burden.
\end{itemize}
\end{frame}

\begin{frame}{Examples and When Not to Use Managed RDBMS}
\begin{block}{Examples}
Amazon RDS, Google Cloud SQL, Azure SQL Database.
\end{block}

\begin{alertblock}{Managed RDBMS may not be suitable if}
\begin{itemize}
  \smallitem
  \item Deep control over the database engine is required.
  \item Data must be stored on premises.
  \item Legal constraints do not allow data to be stored in the cloud.
  \item Cloud costs do not fit the budget.
\end{itemize}
\end{alertblock}
\end{frame}

\subsection{Document Database}

\begin{frame}{Document Database: Concept}
A \kw{document database} stores data as documents, usually in JSON or BSON format.

\vspace{2mm}
Unlike relational databases, document databases do not require every record to have the same structure.
\end{frame}

\begin{frame}[fragile]{Example JSON Document}
\small
\begin{verbatim}
{
  "student_id": 1,
  "full_name": "Nguyen An",
  "major": "Data Science",
  "skills": ["Python", "SQL", "Machine Learning"],
  "address": {
    "city": "Hanoi",
    "district": "Hoan Kiem"
  }
}
\end{verbatim}
\end{frame}

\begin{frame}{When Should We Use a Document Database?}
Use a document database when:
\begin{itemize}
  \smallitem
  \item The data structure is flexible.
  \item Data is naturally represented as JSON.
  \item The schema changes frequently.
  \item The data contains nested structures.
  \item The application evolves quickly and requires flexible data modeling.
\end{itemize}
\end{frame}

\begin{frame}{Suitable Examples and Managed Document Databases}
\begin{block}{Suitable Examples}
User profiles, product catalogs, article content, mobile applications, and content management systems.
\end{block}

\begin{block}{Managed Document Databases}
MongoDB Atlas, Amazon DocumentDB, Couchbase Capella.
\end{block}

A managed document database is useful when fast scaling, replication, sharding, backup, multi-region deployment, access control, and monitoring are needed.
\end{frame}

\subsection{Key-Value Store}

\begin{frame}{Key-Value Store: Concept}
A \kw{key-value database} stores data as pairs of keys and values.

\vspace{2mm}
This database type is simple but provides very fast access speed.
\end{frame}

\begin{frame}[fragile]{Example Key-Value Data}
\small
\begin{verbatim}
user:1001 -> "Nguyen An"
cart:1001 -> ["book", "laptop", "mouse"]
session:abc123 -> "active"
\end{verbatim}
\end{frame}

\begin{frame}{When Should We Use a Key-Value Store?}
Use a key-value store when:
\begin{itemize}
  \smallitem
  \item Extremely fast data access is required.
  \item Data is mainly queried by key.
  \item User sessions need to be stored.
  \item Data caching is required.
  \item Game leaderboards are needed.
  \item Millions of requests per second must be processed.
\end{itemize}
\end{frame}

\begin{frame}{Suitable Examples and Systems}
\begin{block}{Suitable Examples}
Product cache, user sessions, temporary shopping carts, leaderboards, and view counters.
\end{block}

\begin{block}{Example Systems}
Redis, Memcached, Amazon ElastiCache, Azure Cache for Redis, Google Memorystore.
\end{block}
\end{frame}

\subsection{Time-Series Database}

\begin{frame}{Time-Series Database: Concept}
A \kw{time-series database} is designed for data associated with time.

\vspace{2mm}
Each record usually contains:
\begin{itemize}
  \smallitem
  \item A timestamp.
  \item A measured value.
  \item A data source.
\end{itemize}
\end{frame}

\begin{frame}[fragile]{Example Time-Series Data}
\small
\begin{verbatim}
2026-05-31 10:00:00 | sensor_01 | temperature | 32.5
2026-05-31 10:01:00 | sensor_01 | temperature | 32.7
2026-05-31 10:02:00 | sensor_01 | temperature | 32.6
\end{verbatim}
\end{frame}

\begin{frame}{When Should We Use a Time-Series Database?}
Use a time-series database for:
\begin{columns}[T]
\begin{column}{0.48\textwidth}
\begin{itemize}
  \smallitem
  \item Sensor data.
  \item IoT data.
  \item System logs.
  \item Monitoring metrics.
\end{itemize}
\end{column}
\begin{column}{0.48\textwidth}
\begin{itemize}
  \smallitem
  \item Financial data over time.
  \item Machine operation data.
  \item Real-time dashboards.
  \item Historical time-based analysis.
\end{itemize}
\end{column}
\end{columns}
\end{frame}

\begin{frame}{Suitable Examples and Systems}
\begin{block}{Suitable Examples}
Cold-storage temperature monitoring, server CPU/RAM/disk monitoring, truck location tracking over time, stock price analysis, and electricity consumption analysis.
\end{block}

\begin{block}{Example Systems}
InfluxDB Cloud, Amazon Timestream, Azure Data Explorer, TimescaleDB.
\end{block}
\end{frame}

\subsection{Columnar Database}

\begin{frame}{Columnar Database: Concept}
A \kw{columnar database} stores data by columns instead of by rows.

\vspace{2mm}
If the goal is to compute total revenue, the system only needs to read the \texttt{amount} column. Therefore, columnar databases are highly efficient for large-scale data analytics.
\end{frame}

\begin{frame}{Example: Sales Table}
\scriptsize
\begin{center}
\begin{tabular}{llll}
\toprule
\textbf{order\_id} & \textbf{customer\_id} & \textbf{amount} & \textbf{date} \\
\midrule
1 & C01 & 100 & 2026-05-01 \\
2 & C02 & 200 & 2026-05-02 \\
3 & C03 & 150 & 2026-05-03 \\
\bottomrule
\end{tabular}
\end{center}

\vspace{2mm}
For revenue analytics, column-based reading reduces the amount of data that must be scanned.
\end{frame}

\begin{frame}{When Should We Use a Columnar Database?}
Use a columnar database when:
\begin{itemize}
  \smallitem
  \item Large-scale data analysis is required.
  \item BI or data warehousing is needed.
  \item Queries are mainly read and aggregation queries.
  \item Statistical computation over millions or billions of rows is required.
  \item Storage and compute should be separated.
\end{itemize}
\end{frame}

\begin{frame}{Suitable Examples and Systems}
\begin{block}{Suitable Examples}
Revenue reporting, customer behavior analysis, management dashboards, enterprise data warehouses, and logistics data analysis.
\end{block}

\begin{block}{Example Systems}
Google BigQuery, Amazon Redshift, Snowflake, Azure Synapse Analytics.
\end{block}
\end{frame}

\subsection{Graph Database}

\begin{frame}{Graph Database: Concept}
A \kw{graph database} represents data using:
\begin{itemize}
  \smallitem
  \item \kw{Nodes}: entities.
  \item \kw{Edges}: relationships.
  \item \kw{Properties}: attributes.
\end{itemize}

Graph databases are suitable when relationships among objects are more important than individual objects alone.
\end{frame}

\begin{frame}[fragile]{Example Graph Data}
\small
\begin{verbatim}
Minh --friend_of--> An
An --works_at--> Company A
Company A --located_in--> Hanoi
\end{verbatim}
\end{frame}

\begin{frame}{When Should We Use a Graph Database?}
Use a graph database when:
\begin{itemize}
  \smallitem
  \item Data has many complex relationships.
  \item Path or network queries are needed.
  \item Link analysis is required.
  \item Fraud detection is needed.
  \item Recommendation systems need to be built.
\end{itemize}
\end{frame}

\begin{frame}{Suitable Examples and Systems}
\begin{block}{Suitable Examples}
Social networks, friend recommendations, financial fraud analysis, product recommendations, knowledge management, and supply chain dependency analysis.
\end{block}

\begin{block}{Example Systems}
Neo4j Aura, Amazon Neptune, Azure Cosmos DB Gremlin API.
\end{block}
\end{frame}

\section{Database Management and Operations}
\subsection{Managed Database}

\begin{frame}{What Is a Managed Database?}
A \kw{managed database} is a database service managed by a cloud provider on behalf of the technical team.

\vspace{2mm}
The provider may support:
\begin{columns}[T]
\begin{column}{0.48\textwidth}
\begin{itemize}
  \smallitem
  \item Automatic backups.
  \item Data recovery.
  \item Version upgrades.
  \item Resource scaling.
\end{itemize}
\end{column}
\begin{column}{0.48\textwidth}
\begin{itemize}
  \smallitem
  \item Data replication.
  \item Performance monitoring.
  \item Data encryption.
  \item Failover.
\end{itemize}
\end{column}
\end{columns}
\end{frame}

\begin{frame}{Examples of Managed Databases}
\begin{itemize}
  \smallitem
  \item Amazon RDS.
  \item Google Cloud SQL.
  \item Azure Cosmos DB.
  \item MongoDB Atlas.
  \item Amazon Timestream.
  \item Snowflake.
\end{itemize}
\end{frame}

\subsection{When Should We Choose a Managed Database?}
\begin{frame}{When Should We Choose a Managed Database?}
Choose a managed database when:
\begin{itemize}
  \smallitem
  \item The technical team is small.
  \item The team does not want to spend too much time on operations.
  \item High availability is required.
  \item Automatic backup and recovery are needed.
  \item Integrated monitoring is required.
  \item Security and compliance are important.
  \item The team wants to focus on application development rather than infrastructure administration.
\end{itemize}
\end{frame}

\subsection{When Should We Self-Manage a Database?}
\begin{frame}{When Should We Self-Manage a Database?}
Self-management may be suitable when:
\begin{itemize}
  \smallitem
  \item Deep control over system configuration is required.
  \item Special modules must be installed.
  \item Data must be stored on premises.
  \item The organization has a strong DBA team.
  \item Cost optimization is needed in specific cases.
  \item There are special security or legal requirements.
\end{itemize}
\end{frame}

\section{Architectural Patterns Affecting Database Choice}
\subsection{Polyglot Persistence}

\begin{frame}{Polyglot Persistence}
\kw{Polyglot persistence} means that one system uses multiple database types for different needs.

\begin{exampleblock}{E-commerce Example}
\begin{itemize}
  \smallitem
  \item PostgreSQL for orders and payments.
  \item MongoDB for product catalogs.
  \item Redis for caching.
  \item Elasticsearch for search.
  \item Snowflake for data analytics.
\end{itemize}
\end{exampleblock}

The main idea: the entire system does not have to use only one database type.
\end{frame}

\subsection{CQRS}
\begin{frame}{CQRS}
\kw{CQRS} stands for \kw{Command Query Responsibility Segregation}.

\vspace{2mm}
Main idea:
\begin{itemize}
  \smallitem
  \item Separate the write side.
  \item Separate the read side.
\end{itemize}

\begin{exampleblock}{Example}
An SQL database can be used for transactional writes, while a columnar database can be used for reporting and analytics.
\end{exampleblock}

CQRS is suitable when the system has significantly different read and write requirements.
\end{frame}

\subsection{Event Sourcing}
\begin{frame}{Event Sourcing}
\kw{Event sourcing} stores state-changing events instead of storing only the final state.

\begin{block}{Order Example}
\ttfamily OrderCreated \quad PaymentConfirmed \quad OrderPacked \quad OrderShipped \quad OrderDelivered
\end{block}

\begin{itemize}
  \smallitem
  \item Provides a complete change history.
  \item Makes auditing easier.
  \item Allows the system state to be reconstructed at a past point in time.
  \item Fits systems that require traceability.
\end{itemize}
\end{frame}

\subsection{Microservices}
\begin{frame}{Microservices}
In a microservices architecture, each service may manage its own database.

\begin{exampleblock}{Example}
\begin{itemize}
  \smallitem
  \item The user service uses PostgreSQL.
  \item The product service uses MongoDB.
  \item The cache service uses Redis.
  \item The analytics service uses BigQuery.
\end{itemize}
\end{exampleblock}

This helps each service choose the database that best fits its business responsibility.
\end{frame}

\section{Comparison and DBMS Selection Process}
\subsection{Comparison Table}

\begin{frame}{Quick Comparison of Database Types}
\scriptsize
\begin{adjustbox}{max width=\textwidth}
\begin{tabularx}{1.08\textwidth}{>{\raggedright\arraybackslash}p{0.20\textwidth}>{\raggedright\arraybackslash}p{0.19\textwidth}>{\raggedright\arraybackslash}p{0.36\textwidth}>{\raggedright\arraybackslash}X}
\toprule
\textbf{Database Type} & \textbf{Main Data Form} & \textbf{Suitable For} & \textbf{Examples} \\
\midrule
Relational database & Tables, rows, columns & Transactions, structured data, JOINs & MySQL, PostgreSQL \\
Document database & JSON, BSON & Flexible data, changing schema & MongoDB \\
Key-value store & Key-value pairs & Cache, sessions, extremely fast access & Redis \\
Time-series database & Time-based data & IoT, logs, metrics, sensors & InfluxDB \\
Columnar database & Column-based storage & BI, OLAP, data warehouses & BigQuery, Snowflake \\
Graph database & Nodes and edges & Complex relationships, social networks, fraud detection & Neo4j \\
\bottomrule
\end{tabularx}
\end{adjustbox}
\end{frame}

\subsection{Five-Step Selection Process}
\begin{frame}{DBMS Selection Process}
A five-step process can be used:
\begin{enumerate}
  \smallitem
  \item Identify the nature of the data.
  \item Identify the query pattern.
  \item Identify consistency requirements.
  \item Identify scalability requirements.
  \item Identify operational capability.
\end{enumerate}
\end{frame}

\begin{frame}{Step 1: Identify the Nature of the Data}
Questions to answer:
\begin{itemize}
  \smallitem
  \item Does the data have a fixed structure?
  \item Is the data tabular?
  \item Is the data in JSON format?
  \item Is the data time-dependent?
  \item Does the data contain many complex relationships?
\end{itemize}
\end{frame}

\begin{frame}{Step 2: Identify the Query Pattern}
Questions to answer:
\begin{itemize}
  \smallitem
  \item Are many table JOINs needed?
  \item Is fast key-based querying required?
  \item Is large-scale data analysis required?
  \item Are multi-level relationship queries required?
  \item Is a real-time dashboard required?
\end{itemize}
\end{frame}

\begin{frame}{Step 3: Identify Consistency Requirements}
Questions to answer:
\begin{itemize}
  \smallitem
  \item Does the system require ACID transactions?
  \item Can the system accept eventual consistency?
  \item Would temporary data inconsistency cause serious consequences?
\end{itemize}
\end{frame}

\begin{frame}{Step 4: Identify Scalability Requirements}
Questions to answer:
\begin{itemize}
  \smallitem
  \item Does the system need vertical or horizontal scaling?
  \item Is the data growing rapidly?
  \item Can the number of users increase suddenly?
  \item Is multi-region deployment required?
\end{itemize}
\end{frame}

\begin{frame}{Step 5: Identify Operational Capability}
Questions to answer:
\begin{itemize}
  \smallitem
  \item Does the team have database administration experience?
  \item Is there a dedicated DBA?
  \item Is automatic backup required?
  \item Is integrated monitoring required?
  \item Should a managed database be used?
\end{itemize}
\end{frame}

\section{DBMS Selection Scenarios}
\subsection{Student Management System}

\begin{frame}{Scenario 1: Student Management System}
\begin{block}{Characteristics}
\begin{itemize}
  \smallitem
  \item Data has a clear structure.
  \item There are multiple tables such as students, classes, courses, and grades.
  \item Data JOINs are required.
  \item Data integrity must be ensured.
\end{itemize}
\end{block}

\begin{exampleblock}{Suitable Choice}
PostgreSQL or MySQL.
\end{exampleblock}
\end{frame}

\begin{frame}{Scenario 1: Reasoning}
PostgreSQL or MySQL is suitable because:
\begin{itemize}
  \smallitem
  \item They fit relational data.
  \item They support SQL.
  \item They support primary keys and foreign keys.
  \item They make reporting easier.
\end{itemize}
\end{frame}

\subsection{Social Network Application}
\begin{frame}{Scenario 2: Social Network Application}
\begin{block}{Characteristics}
\begin{itemize}
  \smallitem
  \item User data can be flexible.
  \item The system contains posts, comments, and interactions.
  \item There are friendship and following relationships.
  \item Large-scale expansion is required.
\end{itemize}
\end{block}

\begin{exampleblock}{Possible Combined Choice}
MongoDB for content; Redis for cache; a graph database for friendship relationships; a columnar database for analytics.
\end{exampleblock}
\end{frame}

\subsection{IoT Cold-Storage Temperature Monitoring}
\begin{frame}{Scenario 3: IoT Cold-Storage Temperature Monitoring}
\begin{block}{Characteristics}
\begin{itemize}
  \smallitem
  \item Data is continuously written over time.
  \item Each record has a timestamp.
  \item A real-time dashboard is needed.
  \item Historical data storage and aggregation are required.
\end{itemize}
\end{block}

\begin{exampleblock}{Suitable Choice}
InfluxDB or Amazon Timestream.
\end{exampleblock}
\end{frame}

\begin{frame}{Scenario 3: Reasoning}
InfluxDB or Amazon Timestream is suitable because:
\begin{itemize}
  \smallitem
  \item They are optimized for time-series data.
  \item They support retention policies.
  \item They support time-based data analysis.
  \item They fit real-time monitoring dashboards.
\end{itemize}
\end{frame}

\subsection{E-Commerce System}
\begin{frame}{Scenario 4: E-Commerce System}
\begin{block}{Characteristics}
\begin{itemize}
  \smallitem
  \item Orders require high consistency.
  \item Product information can be flexible.
  \item Cache is needed to improve speed.
  \item Revenue analytics is required.
\end{itemize}
\end{block}

\begin{exampleblock}{Combined Choice}
PostgreSQL for orders; MongoDB for product catalogs; Redis for cache; BigQuery or Snowflake for analytics.
\end{exampleblock}
\end{frame}

\section{Implementation Notes and Summary}

\begin{frame}{Implementation Notes}
When deploying a database system, pay attention to:
\begin{itemize}
  \smallitem
  \item Designing backups from the beginning.
  \item Having a disaster recovery plan.
  \item Monitoring query performance.
  \item Monitoring replication errors.
  \item Designing schemas that can evolve.
  \item Avoiding premature optimization before understanding the workload.
  \item Not choosing a database only because it is trendy.
  \item Prioritizing simple solutions if the problem is not yet too complex.
\end{itemize}
\end{frame}

\begin{frame}{Lesson Summary}
In this lesson, we learned that:
\begin{itemize}
  \smallitem
  \item DBMS selection is an important architectural decision.
  \item RDBMSs fit structured data and strong consistency requirements.
  \item Document databases fit flexible JSON-like data.
  \item Key-value stores fit caching and extremely fast access.
  \item Time-series databases fit time-based data.
  \item Columnar databases fit large-scale analytics.
  \item Graph databases fit data with many complex relationships.
  \item Managed databases reduce operational burden.
  \item Large systems may use multiple database types.
\end{itemize}
\end{frame}

\begin{frame}{Conclusion}
\begin{block}{No single database type fits every problem}
A good choice must be based on data characteristics, query patterns, consistency requirements, performance, scalability, operational capability, cost, and the system's long-term direction.
\end{block}

\vspace{2mm}
In modern systems, combining multiple database types is common. The key is to use each database type according to its strengths.
\end{frame}

\begin{frame}{Keywords}
\small
\begin{columns}[T]
\begin{column}{0.33\textwidth}
\begin{itemize}
  \smallitem
  \item DBMS
  \item Relational Database
  \item Document Database
  \item Key-Value Store
  \item Time-Series Database
  \item Columnar Database
  \item Graph Database
\end{itemize}
\end{column}
\begin{column}{0.33\textwidth}
\begin{itemize}
  \smallitem
  \item Managed Database
  \item ACID
  \item Eventual Consistency
  \item Scalability
  \item Vertical Scaling
  \item Horizontal Scaling
  \item Backup
\end{itemize}
\end{column}
\begin{column}{0.33\textwidth}
\begin{itemize}
  \smallitem
  \item Polyglot Persistence
  \item CQRS
  \item Event Sourcing
  \item Microservices
  \item Failover
  \item Workload
  \item Monitoring
\end{itemize}
\end{column}
\end{columns}
\end{frame}

\section{Questions and Exercises}

\begin{frame}[allowframebreaks]{Review Questions}
\small
\begin{enumerate}
  \smallitem
  \item Why should we not choose a DBMS only because it is popular?
  \item Distinguish vertical scaling from horizontal scaling.
  \item Which database type is suitable for a banking system? Why?
  \item When should a document database be used?
  \item In what situations is Redis commonly used?
  \item What type of data is suitable for a time-series database?
  \item How is a columnar database different from a row-oriented database?
  \item What types of problems are suitable for graph databases?
  \item What advantages does a managed database provide?
  \item What is polyglot persistence?
\end{enumerate}
\end{frame}

\begin{frame}[allowframebreaks]{Short Review Quiz}
\small
\textbf{Question 1.} When choosing a DBMS, which factor should be considered first?
\begin{enumerate}
  \qoption{A}{The provider's logo}
  \qoption{B}{The nature of the data and how the system queries it}
  \qoption{C}{Which database is most popular on social media}
  \qoption{D}{The programming language of the user interface}
\end{enumerate}

\textbf{Question 2.} Which database type is usually preferred for a banking system?
\begin{enumerate}
  \qoption{A}{Relational database}
  \qoption{B}{Key-value store}
  \qoption{C}{Graph database}
  \qoption{D}{Time-series database}
\end{enumerate}

\textbf{Question 3.} If data has a flexible JSON structure and the schema changes frequently, what is a suitable choice?
\begin{enumerate}
  \qoption{A}{Document database}
  \qoption{B}{Columnar database}
  \qoption{C}{Time-series database}
  \qoption{D}{Key-value store}
\end{enumerate}

\textbf{Question 4.} Which situation is Redis most suitable for?
\begin{enumerate}
  \qoption{A}{Complex JOIN queries across many tables}
  \qoption{B}{Cache, sessions, or very fast key-based access}
  \qoption{C}{Historical analysis over billions of rows}
  \qoption{D}{Modeling multi-level social relationships}
\end{enumerate}

\textbf{Question 5.} IoT sensor data continuously sent with timestamps should prioritize which database type?
\begin{enumerate}
  \qoption{A}{Graph database}
  \qoption{B}{Time-series database}
  \qoption{C}{Document database}
  \qoption{D}{Relational database}
\end{enumerate}

\textbf{Question 6.} What need is a columnar database most suitable for?
\begin{enumerate}
  \qoption{A}{Storing login sessions}
  \qoption{B}{Analyzing, aggregating, and reporting large-scale data}
  \qoption{C}{Storing friendship relationships in a social network}
  \qoption{D}{Managing ACID banking transactions}
\end{enumerate}

\textbf{Question 7.} What is the main strength of graph databases?
\begin{enumerate}
  \qoption{A}{Querying and analyzing complex relationships among objects}
  \qoption{B}{Storing large image files}
  \qoption{C}{Completely replacing the operating system}
  \qoption{D}{Only storing fixed tabular data}
\end{enumerate}

\textbf{Question 8.} What operational work can a managed database reduce for the technical team?
\begin{enumerate}
  \qoption{A}{Designing the user interface}
  \qoption{B}{Writing all business logic}
  \qoption{C}{Backups, monitoring, version upgrades, and failover}
  \qoption{D}{Eliminating the need for data security}
\end{enumerate}
\end{frame}

\begin{frame}[allowframebreaks]{Application Exercises}
\small
\textbf{Exercise 1.} A startup is building an online order management application. The system includes customer management, product management, order management, payment, revenue reporting, and product search. Propose a suitable database type for each function.

\vspace{3mm}
\textbf{Exercise 2.} A logistics company wants to store truck location data in real time. Each truck sends GPS data every 10 seconds. Answer the following:
\begin{itemize}
  \smallitem
  \item What database type should be used?
  \item Why?
  \item Should a traditional RDBMS be used as the only database?
  \item What additional tool is needed for a real-time dashboard?
\end{itemize}

\vspace{3mm}
\textbf{Exercise 3.} A social networking system needs to store user profiles, posts, comments, friendship relationships, friend recommendations, and view statistics. Propose a data storage architecture that uses multiple database types.
\end{frame}

\begin{frame}{Additional Multiple-Choice Questions}
\small
\textbf{Question 1.} Which database type is most suitable for clearly structured data that requires complex JOINs?
\begin{enumerate}
  \qoption{A}{Key-value store}
  \qoption{B}{Relational database}
  \qoption{C}{Time-series database}
  \qoption{D}{Graph database}
\end{enumerate}

\textbf{Question 2.} What is Redis commonly used for?
\begin{enumerate}
  \qoption{A}{Large-scale data analysis}
  \qoption{B}{Storing complex social relationships}
  \qoption{C}{Cache and sessions}
  \qoption{D}{Column-based data storage}
\end{enumerate}
\end{frame}

\begin{frame}{Additional Multiple-Choice Questions}
\small
\textbf{Question 3.} Sensor data over time is most suitable for which database type?
\begin{enumerate}
  \qoption{A}{Time-series database}
  \qoption{B}{Graph database}
  \qoption{C}{Document database}
  \qoption{D}{Relational database}
\end{enumerate}

\textbf{Question 4.} MongoDB is a popular example of which database type?
\begin{enumerate}
  \qoption{A}{Columnar database}
  \qoption{B}{Document database}
  \qoption{C}{Key-value store}
  \qoption{D}{Time-series database}
\end{enumerate}

\textbf{Question 5.} Graph databases are most suitable for which problem?
\begin{enumerate}
  \qoption{A}{Storing user sessions}
  \qoption{B}{Analyzing complex relationships}
  \qoption{C}{Computing total monthly revenue}
  \qoption{D}{Storing sensor logs}
\end{enumerate}
\end{frame}

\begin{frame}
\centering
\Huge Thank you!\\
\vspace{4mm}
\Large Questions and Discussion
\end{frame}

\end{document}
